\documentclass[10pt]{article}
\setlength\parindent{10pt}
\usepackage{amsmath, amssymb}
\usepackage[margin=1in]{geometry}
\usepackage{mathrsfs}
\usepackage{natbib}
\pagestyle{empty}

\title{Assessing journal and institution influence through reputation recirculation via citations}
\date{Technical Report, \today}
\author{Docampo D. and Cram, L.}
\begin{document}
\maketitle

\begin{abstract}

Traditional citation-based metrics, like the Journal Impact Factor, measure popularity but are vulnerable to manipulation and fail to account for the prestige of the citing source. Eigenvector-based methods (also called spectral ranking) address the desirability of accounting for the influence of citing journals although they remain within the reflexive journal-to-journal citation system and can be manipulated by orchestrating citations.

 We introduce a novel method, the Journal-Institution Reputation Circulation (JIRC) algorithm, that recirculates reputation between journals and institutions via a two-mode citation network to derive more robust influence metrics.
We construct two normalized citation matrices: one from institutions to journals (Matrix A) and the other from journals to institutions (Matrix B). The algorithm solves for the dominant eigenvector of the combined system \(v = P_B^T P_A^T v\), where \(v\) is the institutional influence vector, and derives the journal influence vector \(w\) as \(w = P_A^T v\). We apply JIRC to a large dataset from the ESI Field of Economic \& Business.

The resulting journal and institution classifications show strong construct validity, aligning with established reputations. Moreover, when used to screen a sample of 330 researchers—equally dividing them into a ``Control'' group (Nobel laureates and prestigious award winners) and a ``Test'' group (authors with highly cited papers but no major awards)—JIRC effectively disentangles influence from notoriety. The top tier of researchers was predominantly from the Control group, while the bottom tier was overwhelmingly composed of members from the Test group. The JIRC algorithm thus provides a robust, manipulation-resistant tool for assessing influence at the journal, institutional, and individual researcher level, offering a significant advancement over purely citation-count-based methods.

{\bfseries Keywords}: Scientometrics, Research Evaluation, Citation Analysis, Reputation, Eigenvector Centrality, Pinski-Narin, PageRank, Journal Influence, Institutional Ranking
 \end{abstract}
---

\section{Introduction}

For researchers, review, evaluation and assessment of their work is a high-stakes activity that shapes professional success in publication, recognition, grant funding, appointment, advancement and financial rewards. The temptation to boost, misrepresent, orchestrate or fake research is accordingly high, countered by the opprobrium attached to discovery \cite{mertonSocialStructureScience1996} and moral/ethical values. The balancing of temptation and morality implies that the research literature includes works designed primarily to inflate research performance indicators, rather than reporting findings that advance research. The proportion of `contamination' is not well characterised, but the informed estimate by \cite{jamesheathersApproximately172025} that 1/7 of the literature is fake corresponds to around 700,000 contaminating works per year. Efficient and accurate methods to screen contamination at this level in the research literature would support the systematic, organised scepticism that is so important to scientific progress \cite[Chapt. 9]{zimanRealScienceWhat2001}.

To date, large scale bibliometric studies have been designed generally to identify worthy research or reveal mechanisms of science. Among such approaches, some are applicable to large scale screening for contamination. As discussed by \cite{tahamtanSocialSystemsCitation2022} in their proposal for the Social Systems Citation Theory (SSCT), a bibliometric approach does not deal with humans and their motives but rather begins with and focuses on the communication network of research. Screens that select outliers among research publications using features in bibliographic and reference data have the potential to identify anomalies at scale across much of the research literature. Anomalous works may be ignored, perused through the lens of forensic bibliometrics, or subject to deeper but more expensive textual and peer review. Anomalous work may also help select subjects for psychosocial studies of researchers' individual motives and collective behaviours.

Effective large-scale screening through journal-level bibliometric anomalies has been demonstrated by \cite{kojakuDetectingAnomalousCitation2021}. Their CIDRE (Citation Donors and Recipients) algorithm detects anomalous citation patterns between journals, revealing candidates for possible exclusion from journal indexing services. \cite{chakrabortyIdentificationAnalysisCitation2020} address a similar problem using sets of features to identify anomalies in computer science journals. \cite{evdaimonMetricsDetectSmallScale2024} exhibit author-level screening using the referencing information of tens of millions of articles to quantify the extreme values of field-specific indicator distribution functions that can isolate specific patterns of referencing orchestration. \cite{ioannidisFeaturesSignalsPrecocious2024a, ioannidisEvolvingPatternsExtreme2024a, ioannidisQuantitativeResearchAssessment2023a} identify many indicators that potentially herald large-scale manipulation. \cite{ibrahimCitationManipulationCitation2025} exhibit an author's citation concentration index, $c^2$, the largest number $n$ of papers that cite the author at least $n$ times, showing that it can be used as a warning about manipulation and also revealing how preprint servers may be exploited as additional sources of contamination.

These and similar studies demonstrate that the tails of frequency distributions (FDs) derived from populations of aggregated counts of citations or references may include a higher proportion of contaminating publications. An alternative to examining extrema of FDs is to screen a reference or citation before aggregation, for example by applying a weight.  The field-weighted citation index (FWCI) and relative citation ratio (RCR) are familiar indicators that weight citations to adjust for benign referencing variations between fields of research. The reputable citation (RC) method of \cite{safonScreeningArticlesCitation2025} uses an exogenous rating of citing author's institutions to scale citations prior to aggregation at the journal or author level. Applied to the field of Mathematics, the gap between aggregations of scaled and unscaled citations delineates authors whose elevated citation profile may have little to do with the modest impact and quality of their work.

While the exogenous ranking used in the RC method helps to insulate the indicator from publication contamination, it is challenging to find the large-scale, independent, field-specific institution rankings that the method uses. Accordingly in this paper we use spectral ranking over a multi-partite citation network to establish a scale of prestige that can be conferred on individual references/citations. By comparing raw and scaled indicators, we exhibit a screen that promises to reveal works, authors, institutions and countries that contribute to contamination of the research literature.

The paper is organised as follows. First, we build on reviews of spectral ranking and eigenvalue centrality by \cite{franceschetPageRankStandingShoulders2011, vignaSpectralRanking2016} to clarify distinctions between various measures of relative influence, prestige or status. We present a general approach to eigenvalue centrality over a multi-partite citation network between journals, authors and institutions, and illustrate it using the journal-institution projection of this network's graph. We discuss our findings about use of this approach as a screen, and conclude with a discussion of the limitations and prospects of this family of methods. 

\section{Status and spectral ranking}

The long-standing ISI journal impact factor (JIF) introduced in \cite{garfield1955, garfieldHistoryMeaningJournal2006} reports the average number of citations per article in a journal. The adjustable features of JIF-like indicators include the census and target time intervals and the rules for (a) choosing which articles in the journal are counted in the numerator and denominator of the quotient, and (b) accepting the subset of referening that are counted as citation sources. The JIF takes all citations to have equal value, and thus says nothing about the nature of citing journals. \cite{garfieldHistoryMeaningJournal2006} explains that the JIF helps identify the core of journals in a field, some by large size and others by higher citation frequency (higher impact), and suggests that some authors use the JIF to help select a prestigious journal for publishing their work. \cite{bollenJournalStatus2006a} argues that the JIF is, however, not an indicator of prestige but rather popularity, since it takes no account of the standing of referencing authors. 
 
Many researchers have advocated alternatives to the JIF that do take account of the properties or context of a reference. \cite{waltmanRelationEigenfactorAudience2010} categorise alternatives as (a) normalising the cited-side, (b) normalising the citer-side, (c) using an indicator other than citations/article, or (d) recursively weighting citations. For example, the citer-side adjusted Audience Factor \citep{zitt2008} adjusts the JIF for inter-journal differences by boosting or diminishing the value of a journal's references in proportion to the average number of references per article before summing the weighted citations. This citer-side adjustment compensates for disciplinary differences in the length of reference lists and in the delay before work is cited, both factors that influence the JIF. \cite{habibzadehJournalWeightedImpact2008} propose citer/citing-side scaling that weights a reference by the quotient of the JIFs of citing and cited journal, rescaling this weight to the range 0.1 to 10 over a logit function of the JIF ratio. 

The JIF and many variants focus initially on individual journals characterised by the number, references and citations of their constituent papers. Comparisons that establish the relative standing of journals are made after the JIFs are computed. An alternative view is that journals inhabit a network or social system where relative standing is determined recursively by the whole community. In an early study based on this view, \cite{narinEVALUATIVEBIBLIOMETRICSUSE1976} and \cite{pinski} measure the influence of a publishing unit (journal, department, institution, nation) as the stable distribution of 'influence weight' (a kind of reputation or status) that is exchanged through reference lists. The normalised exchange between $i$ and $j$ is the share of the total number of references given by all units to $j$, $\sum_k C_{k,j}$ that are given by $i$ to $j$, $C_{i,j}$ to. The final influence of each unit is then determined recursively as the influence-weighted sum of the input-output reference ratios. The Index of New Knowledge (INK) proposed by \cite{paulyReinterpretation145influenceWeight1462008} is the initial, non-recursed stage of the Pinski-Narin algorithm.

Pinksi and Narin used the power iteration to determine the principal eigenvector of the input-output matrix, and applied the method to rank Physics Journals. Chapter X of \cite{narinEVALUATIVEBIBLIOMETRICSUSE1976} also presents an application of the method to institutional ranking. The procedure first constructs journal sets and their influence weights (11 fields) and then derives the total influence of around 135 institutions (broadly and by field) by summing the institutional publication count in each journal weighted by the respective influence weight per publication. The algorithmic institutional influences compare favourably with rankings derived from surveys.

Motivated by this work, \cite{gellerCitationInfluenceMethodology1978} considered an alternative exchange of influence or status through the ratio of the number of references from $i$ to $j$ over the total number of references from $i$, $\sum_k C_{ik}$. Geller's normalisation yields a row-stochastic matrix, the transition matrix of a Markov chain. Its principal eigenvector is the long-run probability distribution of influence or status to each unit. \cite{bollenJournalStatus2006a} propose weighted a pageRank algorithm that starts with the same row-normalising citation matrix, and follows \cite{pageMethodNodeRanking1998} by introducing an attenuation factor $0 \le \alpha \le 1$ and a minimum amount of influence (prestige) for each journal. Clarivate's Eigenfactor journal score \citep{bergstromEigenfactorMeasuringValue2007, westEigenfactorMetricsNetwork2010} and the Scimgo-Scopus journal ranking \citep{gonzalez-pereiraNewApproachMetric2010,guerrero-boteFurtherStepForward2012} also use the row-normalised form. 

Journal influence, status or prestige determined by these recursive algorithms refers to the totality of works in the journal. The distribution of per-article influence over a journal set requires normalisation by the number of articles. \cite{waltmanRelationEigenfactorAudience2010} explain these normalisations and show how the per-article audience factors and per-article influence weights can be regarded as variants delineated by the parameter $\alpha$ of the Eigenfactor method's article influence score. They show how audience factors are insensitive to field differences, while influence factors are insensitive to insignificant journals. \citep{vigna} offers an extensive survey of the connections between historical antecedents of recursive algorithms, explaining the important connection between recursions and path summations.

The Reputable Citation algorithm (RC) by Safon, Docampo and Cram \citep{safon2025}  is another recursive algorithm, and the starting point for this note. Recognising that journal influence indicators based on citations alone may be compromised if authors orchestrate reference lists, RC modifies the iterative determination of the relative influence of each journal by injecting exogenous indicators of the prestige of the institutions that publish in it.

\section{Eigenvector based methods}

Consider a set of papers $P$ published in a journal set $J$, each with references referring to other papers and an authorship drawn from sets of institutions $I$ and authors $A$. In this paper we set aside authors and attend to the connections between journals and institutions due to papers. 


The references and authorships can be represented as a table connecting papers $P_{p1}$ and $P_{p2}$ . citation relationship connection the paper to 

We examine in this section the mathematical background of the Pinski \& Narin (PN) principal eigenvector method to later introduce the modifications leading to the new JIRC algorithm.
\subsection{The citation matrix}
Following PN, let's define the ``unit'' as a journal $j$ which is a member of a set of selected
journals $\bf{J}$, where $j = 1 \ldots M$.

Let $C = (c_{ij})$ be a matrix whose elements are the sum of references emitted by
the articles published in journal $i$ (rows) that become citations received by any of the articles published in journal $j$ (columns) during the target
publication window. If we ignore journal self-citations, a common practice, then $c_{ii} = 0$.


\noindent Define the sum of all references from i as
\begin{equation}
s_i = \sum_{j=1}^Mc_{ij}
\end{equation}
The matrix C can be normalised as the row stochastic matrix P where
\begin{equation}
P = (p_{ij}) = (c_{ij}/s_i)
\end{equation}
in which case citations from a journal are normalized by the total number of citations coming from that journal to the other journals in the ensemble.

\textbf{Domingo, this text is correct but I did not previously notice that it is not the PN normalisation. The PN matrix P is unlikely to be stochastic, explaining why Geller (her equation 3) introduces a kind of ``dual'' to the PN matrix to arrive at a Markov interpretation. Vigna p 442 refers to the PN normalisation as ``slightly bizarre''. This helps explain the PN influence-per-article step, as well as the move from the Eigenvector score of a journal to the article influence score or the similar step taken by Bollen etal.  The PN normalisation is to the number of references coming from journal $i$:}

\begin{equation}
P = (p_{ij}) = (c_{ij}/s_j)
\end{equation}

\textbf{This is total number of citations to i - or the total number of references from other units -  divided by total number of references from i -- i.e. input references/ output references. This is not the same as the total number of citations to i divided by the total number of citations from i, which is the recipe for a stochastic matrix}

A stochastic matrix is a transition probability Markov matrix. The repeated product $P^n$ represents a step-wise chain from an initial to a final state
after n steps. If the matrix P is irreducible, there is a single asymptotic final state; if
not, there is no unique final state. It is desirable for computational reasons that the stochastic matrix P  is irreducible. We will assume that in what follows.

\subsection{Influence Factor (Pinski \& Narin - PN)}
Define the journal influence weight vector
$w = (w_j)$ as a column vector of influence weights for each journal in the set. The influence
weight of a journal averages the weighted number of citations from every journal to that
journal so that, ultimately, journals with ``above average'' influence weight are those that
receive more citations from journals with high influence weight and vice versa. This
leads to an implicit equation for w
\begin{equation}
w_j =   \sum_{i=1}^{N}p_{ij}w_i
\end{equation}

\begin{equation}
w = P^Tw
\end{equation}

\textbf{This equation is not the PN equation. Their equation 1 can be written like this one below - I've used the index $j$ on the left, they use $i$}

\begin{equation}
S_j W_j = \sum_{i=1}^{N}C_{ij}
\end{equation}



This equation defines the influence weight for a journal. The influence weight per
publication $w^{\prime}$ is
\begin{equation}
w^{\prime}_j =
\frac{(P^Tw)_j}{a_{j}}
\end{equation}
where $a_j$ is the number of articles published in journal j in the time period
selected.

\noindent PN uses the iteration (called the “Power method”)
\begin{equation}
w(k)=\frac{P^Tw^{(k-1)}}{|P^Tw^{(k-1)}|}
\end{equation}
that converges to the principal eigenvalue and the desired influence vector.

\section{ The new JIRC algorithm}
 The dataset now contains the ensemble of journals,
 $j=1\ldots M$, as well as a collection of institutions, $i=1\ldots N$ with a sizable number of papers  in the field and period under analysis.

When computing journal influence scores, JIRC will inject  values of citing article prestige, $v_i$, based on institutions assessments derived from the  dataset. Conversely, to
 computing institutional influence scores, JIRC will inject  values of citing article prestige, $w_j$, based on journal assessments derived from the same dataset.
We will formulate the problem as a (re)circulation of reputation. The idea is to go from institutions to institutions through journals and vice versa.

\noindent The problem can be stated as follows.

Let's introduce  citation matrix A, from institutions ($i=1\ldots N$) to journals ($j=1\ldots M$), $A = (c_{ij})$,  whose elements are the sum of references emitted by
the articles published by institution $i$ (rows) that become citations received by any of the articles published in journal $j$ (columns) during the target
publication window.

\noindent Let's again define the overall sum of all references from institution $i$ as
\begin{equation}
s_i = \sum_{j=1}^Mc_{ij}
\end{equation}
The matrix A can be normalised as the row stochastic matrix $P_A$, where
\begin{equation}
P_A = (a_{ij}) = (c_{ij}/s_i)
\label{pa}
\end{equation}
in which case references from an institution are normalized by the total number of references from that institution to all journals in the dataset.


We also  introduce  citation matrix B, from journals  to institutions, $B = (c_{ji})$,  whose elements are the sum of references emitted by
the articles published by journal j (rows) that become citations received by any of the articles published in institution i (columns) during the target
publication window.
\begin{equation}
s_j = \sum_{i=1}^Nc_{ji}
\end{equation}
The matrix B can be normalised as the row stochastic matrix $P_B$, where
\begin{equation}
P_B = (b_{ji}) = (c_{ji}/s_j)
\label{pb}
\end{equation}
in which case, references from a journal are normalized by the total number of references from that journal to all the institutions in the dataset.



The influence vectors are connected by the two matrices. $P_A^T$ is $M\times N$, whereas $P_B^T$ is $N\times M$. Both influence column vectors, $w$, $M\times1$, for journals and $v$, $N\times1$, for institutions, receive reputation shots from each other, $w=P_A^Tv$, $v=P_B^Tw$.

Hence, we can now treat the reputation transfer full circle, as
\begin{eqnarray}
 v & = & P_B^TP_A^Tv\\
 w & = & P_A^TP_B^Tw
 \end{eqnarray}
 Matrices $P_A^TP_B^T$ and $P_B^TP_A^T$ are both square, although they have, in general, different sizes: $M\times M$ and $N\times N$, respectively.

 However, it is well-known that they share all the nonzero eigenvalues, the largest in particular, implying that once we find the eigenvector v corresponding to the largest eigenvalue of $P_B^TP_A^T$, we know that the eigenvector corresponding to the largest eigenvalue of $P_A^TP_B^T$ will be $w=P_A^Tv$, and vice versa.

Equations $v=P_B^TP_A^Tv$ and $w=P_A^TP_B^Tw$ define the influence weight for institutions and journals, respectively. The influence weights per
publication, $v^{\prime}$ and $w^{\prime}$ are:
\begin{eqnarray*}
% \nonumber % Remove numbering (before each equation)
v^{\prime}_i &=& \frac{(P_B^TP_A^Tv)_i}{b_{i}} \\
w^{\prime}_j &=& \frac{(P_B^TP_A^Tw)_j}{a_{j}}
\end{eqnarray*}
where $b_i$ and $a_j$ are the number of articles published  in the selected field and time period by institutions and journals, respectively.


\noindent JIRC will also  use the Power method either for institutions or journals, since both will produce the same result:
\begin{equation}
v(k)=\frac{P_B^TP_A^Tv^{(k-1)}}{|P_B^TP_A^Tv^{(k-1)}|}
\label{pm}
\end{equation}
that converges to the eigenvector, $v$, associated with the largest eigenvalue,  and the desired normalized influence vector, $v^{\prime}$. The equation $w=P_A^Tv$ will render the normalized eigenvector for the journal to journal transfer of reputation,
\begin{equation}
w^{\prime}_j =
\frac{(P_A^Tv)_j}{a_{j}}
\end{equation}
\section{Methods}

\subsection{Data Collection and Preparation}
Our analysis is based on publication data from the Web of Science (WoS) Core Collection,  with sizable number of documents  (article and review type) classified within the ESI Field Economics \& Business in the year 2020. This field encompasses four WoS categories: Business, Economics, Finance, and Management.

The dataset includes:
\begin{itemize}
\item Journals (J): $M =613$ journals focused on Economics \& Business.
\item Institutions (I): $N =600$ institutions worldwide with a substantial number of publications in the ESI field Economics \& Business.
\item Published Articles (PA): More than $41,000$ articles published in 2020, classified within the Economics \& Business ESI Field by the Clarivate motor InCites, including journal names and institutional affiliations.
\item Citing Articles (CA): More than $200,000$ articles published between 2020 and 2025  classified within the Economics \& Business ESI Field by the Clarivate motor InCites, citing any of the  articles belonging to PA, including their full citation data and institutional affiliations.
    \end{itemize}
\subsection{Matrix Construction: The JIRC Framework}
As explained before, the core of our method involves the construction of two citation matrices that form a closed loop for reputation recirculation.

\begin{itemize}
\item Matrix $A$ (Institutions → Journals): an (N × M) matrix where each element $a_{ij}$ is the total number of citations from all articles in CA published by institution $i$ to all articles in PA published in journal $j$.
\item Matrix B (Journals → Institutions): an (M × N) matrix where each element \(b_{ji}\) is the total number of citations from all articles in CA published in journal $j$ to all articles published in PA by institution $i$.
    \end{itemize}
Following the Pinski-Narin methodology, we normalize these matrices to be row-stochastic, creating transition matrices \(P_A\) and \(P_B\) (as defined in Eqs. \ref{pa}  and \ref{pb} in the introduction). This normalization controls for the size of institutions and journals, ensuring we measure average influence rather than raw volume.
\subsection{Solving the Reputation Recirculation System}
The influence vectors for institutions, $v$, and journals, $w$, are defined by the coupled system:
\[ v = P_B^T P_A^T v \]
\[ w = P_A^T P_B^T w \]

To check the irreducibility of the two square matrices \(P_B^T P_A^T\) and \(P_A^T P_B^T\), we rely on a standard result in the field of matrix theory, connected to Perron-Frobenius theory, which states that a square matrix is irreducible if and only if its associated directed graph is strongly connected. We  used an algorithm based on \cite{Dijkstra} to settle the issue. Afterwards,
we solve for the principal eigenvector $v$ of the (N × N) matrix \(P_B^T P_A^T\) using the power method (Eq. \ref{pm}). The journal influence vector is then directly calculated as \(w = P_A^T v\).

\subsection{Researcher Screening Application}
To validate the practical utility of JIRC in individual research evaluation, we assembled a list of  330 prominent researchers in Economics \& Business. This list was split into two groups:
\begin{enumerate}
  \item Control Group (C): $165$ Researchers with universally acknowledged influence, including recipients of the Nobel Prize in Economics, the John Bates Clark Medal, Gossen Prize, and other prestigious field-specific awards, as well as European Research Council advanced grant holders.
  \item Test Group (T):$165$ Researchers with a high number of highly-cited papers but without the prestigious awards defining the Control group.
\end{enumerate}
For each paper published by any of the researchers from our sample, we computed a JIRC score as the product of the journals's JIRC score and the best institutional JIRC score among all the institution in the paper byline

Then, for each researcher, we calculated a personal JIRC score by averaging the JIRC scores of all their publications between 2015 and 2024 (matching the  temporal window used by Clarivate in 2025 to identify highly cited papers). Researchers were then ranked and divided into three tiers based on this averaged score.

\section{Results}
\subsection{Journal and Institution Rankings}
The JIRC algorithm produced stable and convergent influence scores for all journals and institutions. The journal rankings ([Table/Appendix 1]) showed strong concordance with expert judgment and other reputation-based rankings, while also introducing nuance by factoring in the prestige of citing institutions. Similarly, the institutional rankings ([Table/Appendix 2]) reflected global academic reputations, effectively identifying centers of genuine intellectual influence.
\subsection{Disentangling Researcher's Influence from Notoriety}
The results of the researcher screening through individual JIRC scores were striking and demonstrate the core strength of the method.
\begin{itemize}
\item Tier 1 (Top 100 Researchers): This group was mainly composed  of researchers from the Control group $(88\%)$. These are individuals whose work is not only widely cited but is cited by influential entities (journals and institutions).
\item Tier 2 (Middle 130 Researchers): This group showed a mix of both Control and Test group members (75 vs 55), indicating a spectrum of influence among highly cited authors.
\item Tier 3 (Bottom 100 Researchers): This group was overwhelmingly populated by researchers from the Test group $(98\%)$. This suggests that while these individuals have achieved high citation counts (notoriety or popularity), their work is cited by sources with lower aggregate influence scores, potentially indicating strategic or circular citation practices within specific sub-communities, or work that is popular but not foundational.
    \end{itemize}
\section{Discussion}

Our proposed JIRC algorithm successfully extends the eigenvector centrality paradigm beyond the journal-journal citation network. By recirculating reputation through a journal-institution loop, we ground journal influence in the prestige of its citing institutions and, conversely, ground institutional prestige in the influence of its citing journals.

The most significant finding is the algorithm's ability to differentiate between influence and mere popularity at the individual researcher level. The clear separation between the award-winning Control group and the highly-cited-but-unawarded Test group across the tiers provides strong empirical validation. It suggests that JIRC is resistant to inflation from tactics like citation stacking or publishing in journals that prioritize rapid publication and high citation volume over intellectual rigor and influence.

The algorithm's accuracy is dependent on clean, disambiguated data for institutions. Future work will explore the impact of self-citations at the institutional and journal level, and the application of JIRC to other scientific fields and to author-level networks directly.
\section{Conclusion}

The Journal-Institution Reputation Circulation algorithm offers a robust and theoretically grounded framework for assessing scholarly influence. It moves beyond counting citations to weighing them based on a recursive definition of reputation that flows between producers (institutions) and disseminators (journals) of knowledge. Our results demonstrate that JIRC is not just a metric for ranking journals and institutions but also a powerful tool for research evaluation, capable of screening individual researchers to identify those with genuine, influential impact as opposed to those who have gained notoriety from a purely citation-based system. In an era of increasing publication volume and potential metric manipulation, JIRC provides a nuanced and reliable measure of true scholarly prestige.


\bibliographystyle{spbasic}
\bibliography{sciomtcs,MyLibrary}

\end{document}

